\centerline{\bf \Huge Принцип крайнего + наглядная геометрия}

\begin{flushright}
        \scriptsize{}
\end{flushright}

\problem{1}
На некоторых двух полях шахматной доски стоят черная и белая пешки. За один ход можно пойти одной из пешек на соседнее поле (две пешки ставить на одно поле нельзя). Можно ли делать такие ходы таким образом, чтобы каждое из возможных расположений пешек встретилось ровно один раз?

\problem{2}
Раскрасьте плоскость в три цвета так, чтобы каждая прямая была окрашена

a) хотя бы в два цвета;

б) ровно в два цвета;

в) ровно в три цвета.

г) Можно ли покрасить её так, чтобы каждая прямая была покрашена ровно в один цвет?

\problem{3}
Может ли во время шахматной партии на каждой из 30 диагоналей оказаться нечётное число фигур?

\problem{4}
Солдаты построены в две шеренги по $n$ человек, так что каждый солдат из первой шеренги не выше стоящего за ним солдата из второй шеренги. В шеренгах солдат выстроили по росту. Докажите, что после этого каждый солдат из первой шеренги также будет не выше стоящего за ним солдата из второй шеренги.

\problem{5}
Плоскость раскрашена в два цвета, причем каждый цвет использован. Верно ли, что найдутся две точки на расстоянии 1 друг от друга

a) разных цветов;

б) одного цвета.


\problem{6}
Разрежьте квадрат на два равных а) пятиугольника б) шестиугольника.

\newpage

\hspace{3mm}

\problem{7}
Если смотреть на аквариум спереди, то рыбка проплыла, как показано на левом рисунке. А если справа --- то как на правом рисунке. Нарисовать вид сверху.

\begin{figure}[h]
    \centering
    \includegraphics[width=0.2\linewidth]{ЭлМШ 2025-2026//VIII блок//7/image.png}
\end{figure}

\problem{8}
Дед звал внука к себе в деревню: "Вот посмотришь, какой я необыкновенный сад посадил! У меня там растут груши и яблони, причём яблони посажены так, что на расстоянии 10 метров от каждой яблони растёт ровно две груши". - "Ну и что тут интересного, - ответил внук. - У тебя, значит, яблонь вдвое меньше, чем груш". "А вот и не угадал, - улыбнулся дед. - Яблонь у меня в саду вдвое больше, чем груш". Нарисуйте, как могли расти яблони и груши в саду у деда.

\problem{9}
На плоскости даны 16 точек (см. рисунок).

a) Покажите, что можно стереть не более восьми из них так, что из оставшихся никакие четыре не будут лежать в вершинах квадрата.

б) Покажите, что можно обойтись стиранием шести точек.

в) Найдите минимальное число точек, которые достаточно стереть для этого.

\begin{figure}[h]
    \centering
    \includegraphics[width=0.2\linewidth]{ЭлМШ 2025-2026//VIII блок//7/image_2.png}
\end{figure}